\documentclass[11pt,a4paper]{article}
\usepackage{brainteaser}

\begin{document}
\brainteasertitle{6}{Les deux œufs}

\begin{statementbox}
Un immeuble possède $100$ étages. Il existe un étage critique $F$, inconnu,
tel qu'un œuf lâché d'un étage supérieur à $F$ se casse, tandis qu'un œuf
lâché de l'étage $F$ ou d'un étage inférieur ne se casse pas.

On dispose de deux œufs identiques. Un œuf intact peut être réutilisé. Quel est
le nombre minimal de lancers nécessaire, dans le pire des cas, pour déterminer
exactement $F$? Donner la stratégie optimale et une preuve de son optimalité.
\end{statementbox}

\begin{hintsbox}
\begin{enumerate}
  \item Avec un seul œuf, pourquoi faut-il tester les étages un par un?
  \item Supposer que l'on veut garantir la réponse en $m$ lancers. À quel
  étage peut-on effectuer le premier lancer?
  \item Si le premier œuf survit, pourquoi le saut suivant doit-il être plus
  petit d'une unité?
  \item Combien d'étages peut-on couvrir avec des sauts $m,m-1,\ldots,1$?
  \item Chercher le plus petit entier $m$ tel que
  \[
  1+2+\cdots+m\geq100.
  \]
\end{enumerate}
\end{hintsbox}

\correction

Supposons que l'on veuille garantir la réponse en au plus $m$ lancers. On
commence à l'étage $m$.

Si l'œuf casse, il reste $m-1$ lancers et un seul œuf pour tester les $m-1$
étages inférieurs un par un. S'il ne casse pas, il ne reste que $m-1$ lancers:
le saut suivant doit être au plus $m-1$, puis $m-2$, etc.

Une stratégie de $m$ lancers peut donc couvrir au maximum
\[
m+(m-1)+\cdots+1=\frac{m(m+1)}2
\]
étages. Il faut
\[
\frac{m(m+1)}2\geq100.
\]
Or
\[
\frac{13\cdot14}{2}=91<100,
\qquad
\frac{14\cdot15}{2}=105\geq100.
\]

Une stratégie optimale teste successivement les étages
\[
14,27,39,50,60,69,77,84,90,95,99,100.
\]
Les écarts théoriques sont $14,13,12,\ldots$; le dernier est tronqué par la
borne $100$. Dès que le premier œuf casse, on reprend étage par étage avec le
second depuis le dernier étage sûr. Le nombre de lancers déjà utilisés plus le
nombre maximal de tests restants ne dépasse jamais $14$.

\begin{resultbox}
\[
\boxed{14\text{ lancers dans le pire des cas}.}
\]
La borne triangulaire prouve simultanément qu'une stratégie en $13$ lancers
est impossible.
\end{resultbox}

\end{document}
